\documentclass[runningheads]{llncs}
\usepackage[T1]{fontenc}
\usepackage{lmodern}
\setlength{\emergencystretch}{2em}
\usepackage{amsmath,amssymb,booktabs,array}
\usepackage{url}
\usepackage[hidelinks]{hyperref}
\begin{document}
\title{Portable Verifiable Policies for Reranking with Partial Source Coverage}
\titlerunning{Portable Verifiable Reranking Policies}
\author{Shenghang Yuan\orcidID{0009-0008-7434-1103}}
\authorrunning{S. Yuan}
\institute{NYU Shanghai, Shanghai, China\\\email{sy4411@nyu.edu}}
\maketitle
\begin{abstract}
Researchers studying constrained stochastic reranking need to pass policies
between a scoring pipeline, an optimizer and an evaluator. When only some
candidates have second-stage scores, that exchange must preserve both missing
coverage and the reference ranking used by the constraint. We present a Python
resource for this workflow: it fixes unsupported positions, exports full-ID
ranking mixtures, and checks worst-case expected-nDCG regret at the recipient.
The checker reconstructs exposure from serialized permutations against a
separately supplied reference, rather than trusting solver-reported exposure.
The resource includes a TREC run adapter, a text-to-policy retrieval example,
a solver-independent checker, downstream evaluation and sampling recipes,
regression tests and an engineering benchmark.
A saved synthetic sweep covers 26 workload/formulation combinations and 78
timed repetitions, with no observed timeout or rejection under a 20-second
per-worker limit and maximum exposure reconstruction discrepancy
$3.46\times10^{-15}$. A bounded historical negative case illustrates why
feasible policies do not by themselves establish useful certificate comparisons.
The contribution is a tested exchange interface for a specific class of
reranking experiments, not a new regret theorem. The evidence establishes
executable integration and numerical consistency; independent human reuse
and retrieval-effectiveness improvement are not established. Code, examples
and aggregate evidence are available in the Apache-2.0 release 0.3.0.
\keywords{Research software \and Stochastic ranking \and Partial coverage \and Policy verification}
\end{abstract}

\section{Introduction}
Consider a researcher who retrieves a candidate list, scores only a subset
with a second-stage model, and sends a randomized reranking policy to a
separate evaluation component. The evaluator needs the same candidate IDs,
coverage mask and reference ranking as the optimizer. Treating absent scores
as zero changes the allowed movement; sending only expected exposure leaves
the evaluator without a distribution to execute. Accepting the reference
embedded in the received policy also permits the sender to change what is
being checked.

Expected exposure is an established way to represent stochastic rankings
\cite{diaz,singh}. This resource addresses a narrow implementation need:
exchange only the positions occupied by scored candidates, optimize a linear
source-score objective, and bound worst-case \emph{expected} nDCG regret
relative to a caller-supplied ranking for arbitrary nonnegative gains.
The output must carry the property through mixture recovery, serialization
and recipient-side checking. Sorting the optimized exposure vector does not
provide that guarantee.

The resource supplies the exchange layer between these components: external
candidate identifiers, coverage-preserving run import, explicit mixtures and
a checker bound to the recipient's own reference. Its core optimizer already
recovers ranking mixtures. The added value is therefore not decomposition
alone, but the handling and checking of the object passed downstream.
The recipient does not trust the solver's claimed exposure or certificate;
it does reuse the same pure-Python mathematical oracle.

We provide a concrete text-to-policy workflow using BM25 and selectively
computed TF-IDF cosine scores. It requires neither relevance labels nor model
downloads and allows replacement of its invented corpus and queries.
A saved engineering sweep characterizes optimization, mixture recovery and
checking separately. A historical negative study is retained as a bounded
use case, not rescued as an effectiveness contribution.

This interface fits experiments that need all three of the following:
partial source coverage, fixed positions for unsupported candidates, and an
expected-regret constraint on an exported stochastic policy. It is not a
drop-in replacement for a system that must emit one deterministic list or
move unscored candidates. Ranker training, relevance evaluation and downstream
sampling remain responsibilities of the surrounding research pipeline.

\section{Relation to Existing Resources}
PyTerrier supports retrieval and reranking pipeline composition and TREC run exchange
\cite{pyterrier}; ranx offers evaluation and fusion methods \cite{ranx}.
Our contribution is not pipeline composition or fusion itself. The additional
interface preserves missing-coverage semantics, exports full ranking mixtures,
and reconstructs exposure against recipient-supplied context to check a
reference-relative expected-regret constraint on the received object.

Exposure-constrained ranking \cite{singh}, stochastic-ranking evaluation
\cite{diaz}, and learning stochastic policies with fairness constraints
\cite{fairpgrank} are established. FairSearch supplies ranked-representation
checks and reranking \cite{fairsearch}, while FairDiverse provides tools for
fairness and diversity experiments \cite{fairdiverse}. Public Birkhoff
implementations already recover permutation mixtures from doubly stochastic
matrices \cite{birkhoff}. The resource builds on these established ideas.

These resources provide adjacent capabilities rather than interchangeable
implementations of the same contract. Our comparison is based on their
documented roles, not an exhaustive claim that no other software can perform
this task. Section~\ref{sec:integration} identifies the concrete responsibilities
handled by this resource and those left to its caller. A reduction in
integration effort remains a question for external reuse, not a conclusion
from the existence of an API.

Statistical risk-control methods answer a different question.
Conformal risk control and its non-exchangeable extensions
\cite{crc,nonexchangeable} depend on calibration observations and assumptions
or explicit distribution-shift terms. Inference-time stochastic ranking with
risk control is a relevant ranking formulation \cite{isrr}.
A pointwise worst-case certificate is not an empirical substitute for those
methods. Conversely, a numerical utility difference does not establish
certificate value when no risk excess was observed to be avoided.
The general idea of a price of robustness is also classical \cite{bertsimas}.

\section{Mathematical Object and Guarantee}
\subsection{Inputs and policy family}
For one query, let $D$ contain $n$ candidates and let $\pi_0$ be the complete
reference permutation. A supported subset $S\subseteq D$ has finite source
scores $s_i$ for exactly its members. Let
$w_1\geq \cdots\geq w_k>0$ be position weights; positions after $k$ have
weight zero. A policy $P$ is a finite distribution over full permutations.
For every component, candidates outside $S$ keep their positions in $\pi_0$.
Thus supported items may exchange their reference slots but cannot displace
unsupported items.

The exposure of item $i$ is
\[
 e_i(P)=\sum_{\pi} P(\pi)w_{\operatorname{pos}_{\pi}(i)},\qquad
 e_i^0=w_{\operatorname{pos}_{\pi_0}(i)}.
\]
The solver maximizes the source-side surrogate
$\sum_{i\in S}s_i e_i(P)$ subject to a regret budget $b\geq0$.
The scores need not be relevance labels. Improving this objective does not
establish improvement in measured nDCG.

For a nonzero nonnegative gain vector $g$, define
\[
 \operatorname{IDCG}_w(g)=\sum_{j=1}^{k} w_j g_{(j)},\qquad
 R(P;g)=\frac{\sum_i g_i(e_i^0-e_i(P))}
                  {\operatorname{IDCG}_w(g)},
\]
where gains are sorted in decreasing order in the denominator.
The all-zero vector contributes zero by convention. The certificate is
$\max(0,\sup_{g\geq0,\,g\ne0} R(P;g))$.
It uses the same weights in actual and ideal gain and is relative to the
specified candidate universe, not documents omitted by first-stage retrieval.

\subsection{Reused top-sum calculation}
Write $\delta_i=e_i^0-e_i(P)$ and let $\delta_{[1]}\geq\cdots\geq
\delta_{[n]}$. The existing oracle computes
\begin{equation}
 R^\star(P)=\max\left(0,\max_{1\leq m\leq n}
 \frac{\sum_{j=1}^{m}\delta_{[j]}}
      {\sum_{j=1}^{\min(m,k)}w_j}\right).
 \label{eq:regret}
\end{equation}
For a binary relevant set of size $m$, the denominator is fixed, and the
largest numerator is the sum of the $m$ largest exposure losses.
For graded gains, a layer-cake decomposition writes $g$ as a nonnegative
combination of nested binary level sets. Both numerator and ideal gain
decompose with those same sets, so their ratio is a weighted average of
binary-set ratios. Binary gains are themselves allowed; hence the bound is
exact for this uncertainty set. We include the reasoning to define what the
software checks, not as a new structural result.

The historical optimizer offers compact and cutting-plane LP formulations
and recovers a mixture of permutations through Birkhoff-style decomposition.
The resource retains these implementations. Equivalent objectives need not
produce identical mixtures, and tied source scores do not promise a
deterministic reference-preserving tie break.

\subsection{Recipient contract and numerical scope}
A common positive rescaling of all position weights cancels in Eq.~1 in real
arithmetic, but not in arbitrary floating-point exposure calculations.
The 0.2.1 interface therefore requires $w_1=1$; callers must normalize custom
weights before constructing their trusted context. Other positive,
non-increasing weights are unchanged. This input restriction is a
correction to 0.2.0: with weight $5\times10^{-324}$, the mixture
$0.6(a,b)+0.4(b,a)$ relative to $(a,b)$ can round to zero reported regret
despite actual regret $0.4$. Version 0.2.0 must not be used at nonunit scales;
the bundled logarithmic-weight workflows already use $w_1=1$. The mathematical
formula remains general; the implementation's supported input domain is narrower.

A policy JSON object contains a schema identifier, context and mixture.
Each mixture component contains a probability and a full ranking of string
IDs. The recipient must supply its own expected context: reference ranking,
supported items, weights and budget. It does not accept the embedded context
as the authority.

Verification rejects unknown fields, duplicate JSON keys, nonfinite numbers,
duplicate or missing candidate IDs, invalid probabilities, altered context and
movement of unsupported items. It checks all components, including those
with zero probability. The total probability mass must differ from one by
at most $10^{-12}$. Accepted probabilities are interpreted after division by
that mass. The checker reconstructs exposure from the rankings and recomputes
Eq.~\ref{eq:regret}; regret must be at most $b+10^{-9}$.
The latter is an explicit floating-point acceptance tolerance, not an
interval-arithmetic proof.

The solver is not imported for checking: the recipient needs only Python's
standard library and the resource's pure-Python oracle. Nevertheless, the
checker does not authenticate the sender, establish optimization optimality,
validate empirical utility, or inspect a live sampler. Replacing both the
policy and the recipient's trusted context is outside its threat model.

The guarantee concerns expectation over the normalized mixture. An individual
sample can incur greater regret. Selecting the largest-probability component
or sorting expected exposure changes the policy and invalidates reuse of its
certificate unless the new object is checked separately.

\section{Software Architecture and Research Workflow}
\label{sec:integration}
\subsection{Interfaces and failure handling}
The distribution separates the optimization core, policy exchange layer,
run-file adapter and engineering harness. The public optimization wrapper maps
external string IDs to core indices, supplies scores on exactly the supported
subset, reconstructs full-ID ranking mixtures, and verifies its own export.
The receiving checker is callable separately on JSON read from disk.

The practical unit of reuse is a query batch with caller-owned reference
contexts, not an LP result file. Table~\ref{tab:integration} shows where the
resource sits in an existing research workflow. A caller can change the source
scorer without changing the exchange format, provided its finite scores and
coverage satisfy the same input contract.

\begin{table}[t]
\caption{Integration responsibilities. The direct core supplies optimization
and mixture recovery; the resource supplies the exchange checks below.}
\label{tab:integration}
\small
\begin{tabular}{p{.22\textwidth}p{.34\textwidth}p{.34\textwidth}}
\toprule
Boundary & Caller retains & Resource handles\\
\midrule
Run import & Reference order, candidate universe and which scores exist &
ID mapping; absent versus zero scores; rejection of extra candidates or queries\\
Policy export & Source scorer, cutoff and budget choice &
Full-ID ranking mixtures; fixed-slot and probability checks after recovery\\
Recipient check & Trusted context stored separately from the received policy &
Exact query-set/context matching; exposure reconstruction; regret recomputation\\
Subsequent use & Evaluation judgments and deployment sampling behavior &
Example consumer: finite-mixture evaluation and sampling diagnostics;
not a live-sampler audit\\
\bottomrule
\end{tabular}
\end{table}

The adapter accepts six-column TREC runs. Explicit ordinal ranks in the target
file determine reference order; source-file scores determine coverage and
objective coefficients. An absent source row means unsupported, not zero.
An observed zero is supported. Extra source queries or documents outside the
target universe are rejected. A query with no source rows receives its
reference policy. Batch verification requires the exact expected query set:
omitting a difficult query is not full-batch success.

Diagnostics include the input role, offending line or query and relevant
numerical discrepancies. The adapter keeps batches in memory and does not
supply a service queue or per-query process isolation. The separate benchmark
harness demonstrates bounded worker execution. Malformed input and solver
failure do not silently become a successful constrained solution.

A direct-solver example makes the wrapper's work inspectable. It manually maps
IDs, calls the same core, maps mixtures back, serializes them and checks against
the original context. Across nine small support/budget settings, source
objectives agree within $10^{-8}$. The comparison checks interface equivalence;
it does not measure user effort. For example, a manual integration must keep
the original string-ID context available after index conversion and must
reject a returned batch that omits a query. The wrapper provides these checks
as explicit interface behavior rather than leaving them in experiment-specific
glue code.

\subsection{From text retrieval to a checked policy}
The supplied workflow uses \texttt{rank-bm25} 0.2.2 and scikit-learn 1.8.0
\cite{rankbm25,sklearn}. Its default corpus contains 20 invented research
snippets and four invented queries, with no qrels or natural-domain labels.
The label-free producer performs the following steps for each query:
\begin{enumerate}
\item Lowercase and tokenize text, retrieve 12 candidates with BM25, breaking
equal scores by document ID.
\item Compute unigram/bigram TF-IDF cosine scores only for the first six
candidates. Features are fitted on the corpus, but second-stage query-document
similarities are evaluated only on that subset.
\item Write target and partial-source TREC files, create requests with cutoff
five and regret budget $0.02$, and optimize policies.
\item Serialize the full mixtures, then reconstruct expected contexts from
the upstream run files and verify the received policies.
\item Write the verification report and a provenance record containing input
hashes, dependency versions, parameters and output hashes.
\end{enumerate}

Users can replace corpus and query files with lists of \texttt{id}/\texttt{text}
objects, or replace the second-stage scoring component while retaining the
run-file interface. IDs must be unique and whitespace-free. The example
refuses to overwrite an existing output directory. A partial directory after
failure is not evidence of successful verification.

Both scorers are inexpensive lexical methods; this is not a claim of neural
reranking or measured computational savings. Its purpose is to exercise a
plausible research boundary using actual scoring libraries rather than
pretyped scores. The default four-query run completed with verified policies.
Regression tests additionally exercise source limits zero, one, six and twelve,
duplicate IDs, extra input fields and malformed text. No relevance-effect
conclusion is drawn from these checks.

\subsection{Using the policy downstream}
A separate, solver-free consumer completes the example through evaluation.
It verifies the received policies against separately retained requests before
computing metrics from evaluation gains. It requires a finite nonnegative gain for every
candidate and rejects missing judgments rather than assigning zero. Gains are
not exponentiated as relevance grades. The ideal denominator uses the same
candidate universe and position weights; all-zero gain queries receive zero
and remain in the reported macro-average.

For each query, the consumer enumerates the normalized mixture to calculate
expected nDCG, cross-checks it against exposure-weighted evaluation, and samples
components with a recorded seed. Sampling supplies diagnostic counts and means,
not the certificate. The text-workflow gain file contains illustrative assignments
added for integration testing, not independent judgments or held-out evidence.
Neither the producer nor the optimizer receives that file. Users can substitute
their own completely judged candidate sets without changing the exchange format.

Five further constructed cases make policy misuse observable. With reference
$[a,f,b]$, supported items $\{a,b\}$, cutoff one and source scores favoring $b$,
budget $0.6$ produces a mixture placing $b$ first with probability $0.6$.
Gains favoring $a$ yield expected nDCG $0.4$; gains favoring $b$ yield $0.6$.
Taking the largest-probability component, or sorting expected exposures within
supported slots, instead puts $b$ first deterministically and has worst-case
regret one. Both replacements fail the original budget. At budget $0.2$, both
replacements retain the reference and pass, yet change expected nDCG from $0.2$
to zero when gains favor $b$. The remaining cases exercise no-coverage fallback
and all-zero gains. These are deliberately chosen counterexamples, not estimates
of failure prevalence. Tests send both the direct-core and public-interface
outputs through this consumer and recover the hand-computable endpoints.

\section{Engineering Evaluation}
\subsection{Questions and evidence}
The evaluation asks whether the exchanged policy preserves its contract, what
the full export-and-check path costs on bounded workloads, and whether it can
be connected to actual scoring libraries. These questions use different
evidence (Table~\ref{tab:evidence}). In particular, running a text pipeline on
invented inputs does not establish natural-domain effectiveness or adoption.

\begin{table}[t]
\caption{Evidence supporting the resource claims. Each row has a distinct
scope; the rows are not independent replications of one scientific effect.}
\label{tab:evidence}
\small
\begin{tabular}{p{.25\textwidth}p{.35\textwidth}p{.30\textwidth}}
\toprule
Evidence & Question addressed & Main limit\\
\midrule
Small exhaustive and interface tests &
Does the implementation match its stated oracle and exchange contract? &
Finite cases; shared oracle, not formal verification\\
26 synthetic workload/formulation combinations &
What are the observed solver, recovery and exchange costs? &
One machine; three repeats; historical 0.1.0 timings\\
BM25/TF-IDF text workflow &
Can real scoring libraries feed the public interfaces? &
Invented text; agent-run, not human reuse\\
Downstream consumer and five constructed cases &
Can a recipient evaluate and sample the policy without silently changing it? &
Illustrative gains; finite checks, not adoption or effectiveness\\
Selected MIND-small aggregates &
Why can a valid certificate comparison still be uninformative? &
Bounded historical case; not an independently reproduced study\\
\bottomrule
\end{tabular}
\end{table}

\subsection{Protocol and evidence boundaries}
The saved synthetic sweep varies candidate count, support size, cutoff,
number of supported positions within the cutoff, score pattern and budget.
It includes empty and singleton support, support below the cutoff, zero
budget, tied scores, partial coverage and full coverage.
There are 26 workload/formulation combinations with three timed repetitions
each. Each worker has a 20-second wall limit and one untimed small warmup.
Errors and timeouts are retained in the reporting structure. All 26
combinations completed in the recorded run.

Measurements were made on Windows 11, an Intel Family 6 Model 197 machine with
14 logical processors, Python 3.12.14, SciPy 1.17.1 and NumPy 2.5.1.
Thread-related environment variables were limited to one, without a claim of
CPU pinning or independently verified solver thread counts.
Three repetitions on one machine support descriptive measurements, not
cross-machine latency guarantees or inferential performance comparisons.

Warm complete time includes request construction, optimization and export
checking, JSON round trip and another verification. Imports, startup and
warmup are recorded separately. LP and mixture-recovery times are instrumented
within this complete path. Peak working set includes interpreter and dependency
overhead. It is not incremental memory allocated only by the LP.

The measurements describe implementation 0.1.0. Preparation 0.1.1 improved
diagnostics and added coverage fields. Public distribution 0.2.0 adds the
retrieval workflow and publication packaging; the optimization core and regret
oracle are unchanged. We do not relabel the historical timings as new-version
measurements.

\subsection{Observed costs and reconstruction}
Table~\ref{tab:timing} gives selected rows; the release's summary script
regenerates the complete table from the saved record without invoking a solver.
Absent measurements remain unavailable rather than zero, and failed statuses
are retained.

\begin{table}[t]
\caption{Selected local measurements: medians of three repetitions.
Matched rows fix support at 100 and active supported slots at five.
Times are milliseconds; memory is peak working set in MiB.}
\label{tab:timing}
\small
\centering
\setlength{\tabcolsep}{4pt}
\begin{tabular}{rrrrrr}
\toprule
$n/|S|$ & Active & Formulation & Complete & Recovery & MiB\\
\midrule
40/20 & 5 & Compact & 2.919 & 0.147 & 69.58\\
40/20 & 5 & Cutting-plane & 31.734 & 0.166 & 70.65\\
200/100 & 5 & Compact, matched & 9.049 & 0.903 & 72.16\\
500/100 & 5 & Compact, matched & 11.333 & 0.931 & 72.62\\
1000/100 & 5 & Compact, matched & 13.435 & 0.946 & 73.00\\
500/500 & 10 & Compact & 125.947 & 33.354 & 97.75\\
\bottomrule
\end{tabular}
\end{table}

The fully covered 500-candidate workload illustrates that mixture recovery is
a material part of the deliverable, not a negligible export detail.
The matched cases hold support size, active supported slots and source scores
fixed. An initial spread-support exploration had reduced active support as
candidate count increased, confounding its scaling interpretation. The matched
comparison was added during engineering exploration; it was not a preregistered
confirmatory design.

Across eight cases solved with both formulations, the maximum discrepancy
between median source objectives is $6.66\times10^{-16}$.
The maximum reconstructed-versus-solver exposure discrepancy across all
repetitions is $3.46\times10^{-15}$. These checks concern numerical consistency,
not relevance or semantic correctness of caller-supplied scores.
The benchmark also exposes an operational limitation: tied source scores
can produce movement without source-objective improvement.

\subsection{Testing and reproducibility}
Tests cover context binding, fixed unsupported slots, strict parsing, invalid
probability mass, regret rejection, empty coverage, query completeness and
round-trip export. Small instances enumerate binary relevance assignments;
graded-gain tests check the reduction and numerical behavior. The release
contains the tests, not just a claimed count.

The CI workflow builds a noneditable wheel and first verifies a sample policy
before installing SciPy. It then installs pinned retrieval/test dependencies,
runs the retrieval and direct-solver examples and executes the portable tests
on Windows and Ubuntu. Both hosted jobs for version 0.3.0 passed 173 tests,
including 22 downstream tests, and completed the text-to-evaluation commands.
A fresh local environment also evaluated those policies using the noneditable
wheel without SciPy. These are installability and interface checks, not
independent human adoption or a production sampler audit.

\section{A Bounded Negative Case and Its Lesson}
The motivating MIND-small study \cite{mind} illustrates the distinction between
policy feasibility and an informative comparison. This case is restricted to
a completed train-only study, six target-blocked ranker pairs, three risk levels
and binary expected-nDCG@10. The same 512 final-evaluation impressions were used
across pair/risk cells; those cells are not independent replications of a domain.

Two implemented statistical controls selected the largest tested path point.
Their measured mean-risk excess supplied no positive excess for the
per-impression certificate to avoid. Utility differences were nevertheless
numerically measurable: the maximum absolute temporal-control-minus-certified
mean difference across 18 cells was approximately $0.000557$.
Thus the problem was not inability to calculate a difference, but lack of
evidence for interpreting that difference as the utility price of \emph{observed
avoided risk}. This does not establish policy equivalence, rule out ex-ante
assurance value, or imply the adversarial constraint was slack.

The release supplies selected aggregate decision rows with the source aggregate
hash. This is a transparent secondary record, not a fully independently
reproducible natural-domain study: raw data, individual effects and private
execution receipts are not redistributed. The historical case is not the
resource's effectiveness claim. The earlier held-out pilot lacks a recoverable
effect distribution and is not pooled with this case or made a headline result.

The transferable experimental lesson is modest: check whether the intended
comparison is active before scaling a certificate study. A valid exported
policy is necessary for that study, but cannot make an inactive empirical
comparison scientifically informative. Neither a zero observed mean-risk excess
nor this small sample establishes a population per-impression violation bound.

\section{Availability, Maintenance and Limitations}
Version 0.3.0 is distributed at:
\begin{center}\small
\url{https://github.com/sy4411-sudo/partial-coverage-ranking-resource}
\end{center}
The tagged release contains a source archive, wheel, an author-manuscript
PDF and SHA-256 checksums. Version 0.3.0 includes the downstream evaluation
examples and this revised manuscript, retaining the numerical-input correction
introduced in 0.2.1. The earlier 0.2.0 and 0.2.1 release assets remain unchanged.
The source inventory records the sizes and hashes of
allowlisted payload files. GitHub releases provide versioned access, but are
not a cryptographically immutable preservation service. No DOI or PyPI
publication is claimed. The package has no mandatory solver dependency;
the retrieval extra pins the tested scoring stack.

Code, documentation and invented examples use Apache-2.0. Third-party
dependencies retain their licenses, and no MIND data or private feedback is
redistributed. Citation metadata names the author, affiliation and ORCID.
Shenghang Yuan is the maintainer; the public repository's issue tracker is the
support channel. Reports should include the release, Python/dependency versions,
platform and a minimal synthetic example, not private identifiers or raw logs.
The maintenance plan prioritizes verifier defects and broken documented
workflows, adds regression tests for confirmed defects, and reruns installed-wheel
CI before a corrective release. Fixes receive a new version rather than replacing
0.2.0 assets. Changes to schemas or acceptance semantics must be documented;
an old policy is not silently reinterpreted under a changed certificate contract.
There is no production SLA or guaranteed maintenance duration.

The declared Python minimum is 3.11, whereas the demonstrated release CI matrix
is Python 3.12 on Windows and Ubuntu. The latter, with the pinned optional
dependencies, is the recommended reproduction environment. The declaration is
not evidence that every later Python version, macOS or alternative dependency
stack has been tested.

Coverage and reference quality are also limitations. Fixing unsupported slots
can preserve disparities already present in the reference ranking, and an
uneven source-coverage pattern limits which candidates can benefit from
reranking. The certificate neither corrects candidate omissions nor certifies
group fairness, exposure parity or counterfactual equity. Users must assess
coverage and reference bias separately for their own application; no such
fairness evaluation is claimed here.

Several limitations are material. The recipient and solver share a mathematical
oracle; finite tests are not a formal proof of every implementation path.
The API handles in-memory batches, not distributed production execution.
There is no live-sampler audit. The lexical example demonstrates integration,
not effectiveness on a second natural domain. Timing evidence is local and
historical. Independent human reuse and measured integration effort remain
unavailable. The manuscript is an unsubmitted author draft, not a claim of
acceptance or complete adoption evidence.

\section{Conclusion}
This resource makes the boundary between constrained ranking optimization
and exchanged stochastic policies explicit. It supplies fixed-slot partial
coverage, full-ID mixtures, recipient-bound verification, TREC interoperability,
a text-to-policy integration and inspectable engineering evidence.
The downstream example connects that output to explicit evaluation-only gains,
finite-mixture metrics and sampling diagnostics, including cases where
deterministic substitutions change or invalidate the policy.
It is intended for researchers who need to transfer and check this particular
policy object without reimplementing its exchange semantics. The tests,
measured workloads and text integration support that technical capability;
whether it helps another researcher complete their own integration remains
the next evidence gap. The negative case reinforces a separate lesson:
verify that a comparison actually activates before scaling the experiment.

\paragraph{AI assistance disclosure.}
AI tools assisted coding, checks and drafting; the author retains responsibility.
Automated checks are not independent human validation.

\clearpage
\begin{thebibliography}{15}
\bibitem{diaz} Diaz, F., Mitra, B., Ekstrand, M.D., Biega, A.J., Carterette, B.:
Evaluating Stochastic Rankings with Expected Exposure. SIGIR, 113--122 (2020).
\doi{10.1145/3397271.3401098}
\bibitem{singh} Singh, A., Joachims, T.: Fairness of Exposure in Rankings.
KDD, 2219--2228 (2018). \doi{10.1145/3219819.3220088}
\bibitem{pyterrier} PyTerrier contributors: Operators and I/O documentation.
\url{https://pyterrier.readthedocs.io/en/stable/operators.html};
\url{https://pyterrier.readthedocs.io/en/stable/io.html}. Accessed 7 October 2026.
\bibitem{ranx} ranx contributors: Fusion documentation.
\url{https://amenra.github.io/ranx/fusion/}. Accessed 7 October 2026.
\bibitem{fairpgrank} Singh, A., Joachims, T.: Policy Learning for Fairness in Ranking.
NeurIPS (2019). \url{https://github.com/ashudeep/Fair-PGRank}
\bibitem{fairsearch} FairSearch contributors: FairSearch Python library.
\url{https://github.com/fair-search/fairsearch-fair-python}. Accessed 7 October 2026.
\bibitem{fairdiverse} FairDiverse contributors: FairDiverse research platform.
\url{https://github.com/XuChen0427/FairDiverse}. Accessed 7 October 2026.
\bibitem{birkhoff} Finkelstein, J.: Birkhoff decomposition implementation.
\url{https://github.com/jfinkels/birkhoff}. Accessed 7 October 2026.
\bibitem{crc} Angelopoulos, A.N., Bates, S., Fisch, A., Lei, L., Schuster, T.:
Conformal Risk Control. ICLR (2024). \url{https://arxiv.org/abs/2208.02814}
\bibitem{nonexchangeable} Farinhas, A., Zerva, C., Ulmer, D., Martins, A.F.T.:
Non-Exchangeable Conformal Risk Control. ICLR (2024).
\url{https://arxiv.org/abs/2310.01262}
\bibitem{isrr} Guo, R., Ton, J.-F., Liu, Y., Li, H.:
Inference-time Stochastic Ranking with Risk Control (2023).
\url{https://arxiv.org/abs/2306.07188}
\bibitem{bertsimas} Bertsimas, D., Sim, M.: The Price of Robustness.
Operations Research 52(1), 35--53 (2004). \doi{10.1287/opre.1030.0065}
\bibitem{rankbm25} Brown, D.: Rank-BM25, version 0.2.2.
\url{https://github.com/dorianbrown/rank_bm25}. Accessed 7 October 2026.
\bibitem{sklearn} scikit-learn contributors: TfidfVectorizer documentation.
\url{https://scikit-learn.org/stable/modules/generated/sklearn.feature_extraction.text.TfidfVectorizer.html}.
Accessed 7 October 2026.
\bibitem{mind} Wu, F., et al.: MIND: A Large-scale Dataset for News Recommendation.
ACL, 3597--3606 (2020). \doi{10.18653/v1/2020.acl-main.331}
\end{thebibliography}
\end{document}
